% --- Archon physics formalization source begin ---
% archon:physics
% archon:covers PhyXMiniProblems/problem_phyx_mini_0779.lean
% archon:source-report reports/phyx_mini/problem_phyx_mini_0779.source.json

\chapter{Physics problem phyx\_mini\_0779}
\label{ch:PhyXMiniProblems_problem_phyx_mini_0779}

\paragraph{Problem source.}
\#\# Physical scenario

A string is wrapped several times around the rim of a small hoop with radius \textbackslash{}(8.00\textbackslash{}text\{ cm\}\textbackslash{}) and mass \textbackslash{}(0.180\textbackslash{}text\{ kg\}\textbackslash{}). The free end of the string is held in place and the hoop is released from rest as shown in figure. After the hoop has descended \textbackslash{}(75.0\textbackslash{}text\{ cm\}\textbackslash{}).

\#\# Question

calculate the speed of its center.

\#\# Answer choices

- A: A: 3.84m/s

- B: B: 2.71m/s

- C: C: 1.56m/s

- D: D: 3.21m/s

\#\# Figure caption (auxiliary; use the image as primary evidence)

The image shows a hand holding a vertical string wrapped around the edge of a cylindrical object. Arrows on the cylinder indicate rotational movement, suggesting the cylinder can spin around its vertical axis. The text "0.0800 m" is present, likely indicating the radius or diameter of the cylinder.

\#\# Dataset metadata

Rotational Motion; Physical Model Grounding Reasoning, Multi-Formula Reasoning

\paragraph{Recorded answer/context.}
B: B: 2.71m/s

\paragraph{Figure/image path.}
/root/proposal\_for\_physic/science-mango/phyx\_mini\_run/phyx\_data/test\_image/779.png

\paragraph{Formalization target.}
create a compiling Lean file with sorry bodies at `PhyXMiniProblems/problem\_phyx\_mini\_0779.lean`. The Lean declarations must preserve the physical quantities, dimensions or dimensional roles, figure labels, governing-law hypotheses, and final relation expressed by this problem.
Use Mathlib/Physlib names found through LeanExplore where available. If a physics API is missing, introduce faithful local abstractions rather than scalar placeholder aliases.

\begin{theorem}[Physics formalization target]
\label{thm:physics:phyx_mini_0779:target}
\lean{PhyXMiniProblems.ProblemPhyXMini0779.problem_phyx_mini_0779}
\uses{def:physics:phyx-mini-0779:phyxminiproblems-problemphyxmini0779-speedinmeterspersecond, def:physics:phyx-mini-0779:phyxminiproblems-problemphyxmini0779-unwindinghoopsetup, def:physics:phyx-mini-0779:phyxminiproblems-problemphyxmini0779-matchesproblemscenario, def:physics:phyx-mini-0779:phyxminiproblems-problemphyxmini0779-matchesproblemreadouts, def:physics:phyx-mini-0779:phyxminiproblems-problemphyxmini0779-matchesprimaryfigure, def:physics:phyx-mini-0779:phyxminiproblems-problemphyxmini0779-usesstandardearthgravity, def:physics:phyx-mini-0779:phyxminiproblems-problemphyxmini0779-hasphysicalparameters, def:physics:phyx-mini-0779:phyxminiproblems-problemphyxmini0779-satisfiesthinhoopinertialaw, def:physics:phyx-mini-0779:phyxminiproblems-problemphyxmini0779-satisfiesfixedstringnoslipkinematics, def:physics:phyx-mini-0779:phyxminiproblems-problemphyxmini0779-conservesmechanicalenergyduringdescent, def:physics:phyx-mini-0779:phyxminiproblems-problemphyxmini0779-roundstodisplayedhundredth, def:physics:phyx-mini-0779:phyxminiproblems-problemphyxmini0779-isuniqueclosestdisplayedspeed}
The assigned autoformalize agent should translate this physics problem into Lean declarations in the covered file, with theorem and lemma proof bodies written as `by sorry`.
\end{theorem}
\begin{proof}
This is an autoformalization task, not a proof task. Produce statements that can later be proved by the physics prover without weakening the physical model.
\end{proof}

% --- Archon declaration topology begin ---
\section{Declaration topology}
\begin{definition}[acceleration Dimension]
  \label{def:physics:phyx-mini-0779:phyxminiproblems-problemphyxmini0779-accelerationdimension}
  \lean{PhyXMiniProblems.ProblemPhyXMini0779.accelerationDimension}
  The physical dimension L T⁻² of an acceleration magnitude
\end{definition}
\begin{proof}
  This is the declared model definition of acceleration Dimension.
\end{proof}
\begin{definition}[moment Of Inertia Dimension]
  \label{def:physics:phyx-mini-0779:phyxminiproblems-problemphyxmini0779-momentofinertiadimension}
  \lean{PhyXMiniProblems.ProblemPhyXMini0779.momentOfInertiaDimension}
  The physical dimension M L² of an axial moment of inertia
\end{definition}
\begin{proof}
  This is the declared model definition of moment Of Inertia Dimension.
\end{proof}
\begin{definition}[Mass Quantity]
  \label{def:physics:phyx-mini-0779:phyxminiproblems-problemphyxmini0779-massquantity}
  \lean{PhyXMiniProblems.ProblemPhyXMini0779.MassQuantity}
  A nonnegative, unit-independent physical mass
\end{definition}
\begin{proof}
  This is the declared model definition of Mass Quantity.
\end{proof}
\begin{definition}[Length Quantity]
  \label{def:physics:phyx-mini-0779:phyxminiproblems-problemphyxmini0779-lengthquantity}
  \lean{PhyXMiniProblems.ProblemPhyXMini0779.LengthQuantity}
  A nonnegative, unit-independent physical length
\end{definition}
\begin{proof}
  This is the declared model definition of Length Quantity.
\end{proof}
\begin{definition}[Acceleration Quantity]
  \label{def:physics:phyx-mini-0779:phyxminiproblems-problemphyxmini0779-accelerationquantity}
  \lean{PhyXMiniProblems.ProblemPhyXMini0779.AccelerationQuantity}
  \uses{def:physics:phyx-mini-0779:phyxminiproblems-problemphyxmini0779-accelerationdimension}
  A nonnegative, unit-independent acceleration magnitude
\end{definition}
\begin{proof}
  This is the declared model definition of Acceleration Quantity.
\end{proof}
\begin{definition}[Angular Speed Quantity]
  \label{def:physics:phyx-mini-0779:phyxminiproblems-problemphyxmini0779-angularspeedquantity}
  \lean{PhyXMiniProblems.ProblemPhyXMini0779.AngularSpeedQuantity}
  A nonnegative angular speed; radians are dimensionless
\end{definition}
\begin{proof}
  This is the declared model definition of Angular Speed Quantity.
\end{proof}
\begin{definition}[Moment Of Inertia Quantity]
  \label{def:physics:phyx-mini-0779:phyxminiproblems-problemphyxmini0779-momentofinertiaquantity}
  \lean{PhyXMiniProblems.ProblemPhyXMini0779.MomentOfInertiaQuantity}
  \uses{def:physics:phyx-mini-0779:phyxminiproblems-problemphyxmini0779-momentofinertiadimension}
  A nonnegative axial moment of inertia
\end{definition}
\begin{proof}
  This is the declared model definition of Moment Of Inertia Quantity.
\end{proof}
\begin{definition}[mass Readout]
  \label{def:physics:phyx-mini-0779:phyxminiproblems-problemphyxmini0779-massreadout}
  \lean{PhyXMiniProblems.ProblemPhyXMini0779.massReadout}
  \uses{def:physics:phyx-mini-0779:phyxminiproblems-problemphyxmini0779-massquantity}
  Read a mass in a selected Physlib mass unit
\end{definition}
\begin{proof}
  This is the declared model definition of mass Readout.
\end{proof}
\begin{definition}[length Readout]
  \label{def:physics:phyx-mini-0779:phyxminiproblems-problemphyxmini0779-lengthreadout}
  \lean{PhyXMiniProblems.ProblemPhyXMini0779.lengthReadout}
  \uses{def:physics:phyx-mini-0779:phyxminiproblems-problemphyxmini0779-lengthquantity}
  Read a length in a selected Physlib length unit
\end{definition}
\begin{proof}
  This is the declared model definition of length Readout.
\end{proof}
\begin{definition}[speed Readout]
  \label{def:physics:phyx-mini-0779:phyxminiproblems-problemphyxmini0779-speedreadout}
  \lean{PhyXMiniProblems.ProblemPhyXMini0779.speedReadout}
  Read a speed in coherent selected length and time units
\end{definition}
\begin{proof}
  This is the declared model definition of speed Readout.
\end{proof}
\begin{definition}[acceleration Readout]
  \label{def:physics:phyx-mini-0779:phyxminiproblems-problemphyxmini0779-accelerationreadout}
  \lean{PhyXMiniProblems.ProblemPhyXMini0779.accelerationReadout}
  \uses{def:physics:phyx-mini-0779:phyxminiproblems-problemphyxmini0779-accelerationquantity}
  Read an acceleration in coherent selected length and time units
\end{definition}
\begin{proof}
  This is the declared model definition of acceleration Readout.
\end{proof}
\begin{definition}[angular Speed Readout]
  \label{def:physics:phyx-mini-0779:phyxminiproblems-problemphyxmini0779-angularspeedreadout}
  \lean{PhyXMiniProblems.ProblemPhyXMini0779.angularSpeedReadout}
  \uses{def:physics:phyx-mini-0779:phyxminiproblems-problemphyxmini0779-angularspeedquantity}
  Read angular speed in radians per selected time unit
\end{definition}
\begin{proof}
  This is the declared model definition of angular Speed Readout.
\end{proof}
\begin{definition}[moment Of Inertia Readout]
  \label{def:physics:phyx-mini-0779:phyxminiproblems-problemphyxmini0779-momentofinertiareadout}
  \lean{PhyXMiniProblems.ProblemPhyXMini0779.momentOfInertiaReadout}
  \uses{def:physics:phyx-mini-0779:phyxminiproblems-problemphyxmini0779-momentofinertiaquantity}
  Read a moment of inertia in coherent selected mass and length units
\end{definition}
\begin{proof}
  This is the declared model definition of moment Of Inertia Readout.
\end{proof}
\begin{definition}[mass In Kilograms]
  \label{def:physics:phyx-mini-0779:phyxminiproblems-problemphyxmini0779-massinkilograms}
  \lean{PhyXMiniProblems.ProblemPhyXMini0779.massInKilograms}
  \uses{def:physics:phyx-mini-0779:phyxminiproblems-problemphyxmini0779-massquantity, def:physics:phyx-mini-0779:phyxminiproblems-problemphyxmini0779-massreadout}
  Kilogram readout of a physical mass
\end{definition}
\begin{proof}
  This is the declared model definition of mass In Kilograms.
\end{proof}
\begin{definition}[length In Meters]
  \label{def:physics:phyx-mini-0779:phyxminiproblems-problemphyxmini0779-lengthinmeters}
  \lean{PhyXMiniProblems.ProblemPhyXMini0779.lengthInMeters}
  \uses{def:physics:phyx-mini-0779:phyxminiproblems-problemphyxmini0779-lengthquantity, def:physics:phyx-mini-0779:phyxminiproblems-problemphyxmini0779-lengthreadout}
  Meter readout of a physical length
\end{definition}
\begin{proof}
  This is the declared model definition of length In Meters.
\end{proof}
\begin{definition}[length In Centimeters]
  \label{def:physics:phyx-mini-0779:phyxminiproblems-problemphyxmini0779-lengthincentimeters}
  \lean{PhyXMiniProblems.ProblemPhyXMini0779.lengthInCentimeters}
  \uses{def:physics:phyx-mini-0779:phyxminiproblems-problemphyxmini0779-lengthquantity, def:physics:phyx-mini-0779:phyxminiproblems-problemphyxmini0779-lengthreadout}
  Centimeter readout used by the problem statement
\end{definition}
\begin{proof}
  This is the declared model definition of length In Centimeters.
\end{proof}
\begin{definition}[speed In Meters Per Second]
  \label{def:physics:phyx-mini-0779:phyxminiproblems-problemphyxmini0779-speedinmeterspersecond}
  \lean{PhyXMiniProblems.ProblemPhyXMini0779.speedInMetersPerSecond}
  \uses{def:physics:phyx-mini-0779:phyxminiproblems-problemphyxmini0779-speedreadout}
  SI readout of a speed in meters per second
\end{definition}
\begin{proof}
  This is the declared model definition of speed In Meters Per Second.
\end{proof}
\begin{definition}[acceleration In Meters Per Second Squared]
  \label{def:physics:phyx-mini-0779:phyxminiproblems-problemphyxmini0779-accelerationinmeterspersecondsquared}
  \lean{PhyXMiniProblems.ProblemPhyXMini0779.accelerationInMetersPerSecondSquared}
  \uses{def:physics:phyx-mini-0779:phyxminiproblems-problemphyxmini0779-accelerationquantity, def:physics:phyx-mini-0779:phyxminiproblems-problemphyxmini0779-accelerationreadout}
  SI readout of an acceleration in meters per second squared
\end{definition}
\begin{proof}
  This is the declared model definition of acceleration In Meters Per Second Squared.
\end{proof}
\begin{definition}[angular Speed In Radians Per Second]
  \label{def:physics:phyx-mini-0779:phyxminiproblems-problemphyxmini0779-angularspeedinradianspersecond}
  \lean{PhyXMiniProblems.ProblemPhyXMini0779.angularSpeedInRadiansPerSecond}
  \uses{def:physics:phyx-mini-0779:phyxminiproblems-problemphyxmini0779-angularspeedquantity, def:physics:phyx-mini-0779:phyxminiproblems-problemphyxmini0779-angularspeedreadout}
  SI readout of an angular speed in radians per second
\end{definition}
\begin{proof}
  This is the declared model definition of angular Speed In Radians Per Second.
\end{proof}
\begin{definition}[moment Of Inertia In Kilogram Meters Squared]
  \label{def:physics:phyx-mini-0779:phyxminiproblems-problemphyxmini0779-momentofinertiainkilogrammeterssquared}
  \lean{PhyXMiniProblems.ProblemPhyXMini0779.momentOfInertiaInKilogramMetersSquared}
  \uses{def:physics:phyx-mini-0779:phyxminiproblems-problemphyxmini0779-momentofinertiaquantity, def:physics:phyx-mini-0779:phyxminiproblems-problemphyxmini0779-momentofinertiareadout}
  SI readout of a moment of inertia in kilogram-meter squared
\end{definition}
\begin{proof}
  This is the declared model definition of moment Of Inertia In Kilogram Meters Squared.
\end{proof}
\begin{definition}[Hoop Mass Model]
  \label{def:physics:phyx-mini-0779:phyxminiproblems-problemphyxmini0779-hoopmassmodel}
  \lean{PhyXMiniProblems.ProblemPhyXMini0779.HoopMassModel}
  Idealized mass distribution associated with the word "hoop"
\end{definition}
\begin{proof}
  This is the declared model definition of Hoop Mass Model.
\end{proof}
\begin{definition}[String Model]
  \label{def:physics:phyx-mini-0779:phyxminiproblems-problemphyxmini0779-stringmodel}
  \lean{PhyXMiniProblems.ProblemPhyXMini0779.StringModel}
  Idealization of the string whose own energy is neglected
\end{definition}
\begin{proof}
  This is the declared model definition of String Model.
\end{proof}
\begin{definition}[String Hoop Contact]
  \label{def:physics:phyx-mini-0779:phyxminiproblems-problemphyxmini0779-stringhoopcontact}
  \lean{PhyXMiniProblems.ProblemPhyXMini0779.StringHoopContact}
  Contact condition between the wound string and the rim
\end{definition}
\begin{proof}
  This is the declared model definition of String Hoop Contact.
\end{proof}
\begin{definition}[String End Condition]
  \label{def:physics:phyx-mini-0779:phyxminiproblems-problemphyxmini0779-stringendcondition}
  \lean{PhyXMiniProblems.ProblemPhyXMini0779.StringEndCondition}
  Constraint on the upper free end of the string
\end{definition}
\begin{proof}
  This is the declared model definition of String End Condition.
\end{proof}
\begin{definition}[Center Motion Direction]
  \label{def:physics:phyx-mini-0779:phyxminiproblems-problemphyxmini0779-centermotiondirection}
  \lean{PhyXMiniProblems.ProblemPhyXMini0779.CenterMotionDirection}
  Direction of the center-of-mass motion in the pictured experiment
\end{definition}
\begin{proof}
  This is the declared model definition of Center Motion Direction.
\end{proof}
\begin{definition}[Rotation Sense]
  \label{def:physics:phyx-mini-0779:phyxminiproblems-problemphyxmini0779-rotationsense}
  \lean{PhyXMiniProblems.ProblemPhyXMini0779.RotationSense}
  Rotation sense as viewed in the plane of the supplied image
\end{definition}
\begin{proof}
  This is the declared model definition of Rotation Sense.
\end{proof}
\begin{definition}[Figure Object]
  \label{def:physics:phyx-mini-0779:phyxminiproblems-problemphyxmini0779-figureobject}
  \lean{PhyXMiniProblems.ProblemPhyXMini0779.FigureObject}
  Individually visible objects or annotations in image 779.png
\end{definition}
\begin{proof}
  This is the declared model definition of Figure Object.
\end{proof}
\begin{definition}[Supplied Hoop Figure]
  \label{def:physics:phyx-mini-0779:phyxminiproblems-problemphyxmini0779-suppliedhoopfigure}
  \lean{PhyXMiniProblems.ProblemPhyXMini0779.SuppliedHoopFigure}
  \uses{def:physics:phyx-mini-0779:phyxminiproblems-problemphyxmini0779-rotationsense, def:physics:phyx-mini-0779:phyxminiproblems-problemphyxmini0779-figureobject}
  Structured transcription of the primary raster
\end{definition}
\begin{proof}
  This is the declared model definition of Supplied Hoop Figure.
\end{proof}
\begin{definition}[Unwinding Hoop Setup]
  \label{def:physics:phyx-mini-0779:phyxminiproblems-problemphyxmini0779-unwindinghoopsetup}
  \lean{PhyXMiniProblems.ProblemPhyXMini0779.UnwindingHoopSetup}
  \uses{def:physics:phyx-mini-0779:phyxminiproblems-problemphyxmini0779-massquantity, def:physics:phyx-mini-0779:phyxminiproblems-problemphyxmini0779-lengthquantity, def:physics:phyx-mini-0779:phyxminiproblems-problemphyxmini0779-accelerationquantity, def:physics:phyx-mini-0779:phyxminiproblems-problemphyxmini0779-angularspeedquantity, def:physics:phyx-mini-0779:phyxminiproblems-problemphyxmini0779-momentofinertiaquantity, def:physics:phyx-mini-0779:phyxminiproblems-problemphyxmini0779-hoopmassmodel, def:physics:phyx-mini-0779:phyxminiproblems-problemphyxmini0779-stringmodel, def:physics:phyx-mini-0779:phyxminiproblems-problemphyxmini0779-stringhoopcontact, def:physics:phyx-mini-0779:phyxminiproblems-problemphyxmini0779-stringendcondition, def:physics:phyx-mini-0779:phyxminiproblems-problemphyxmini0779-centermotiondirection, def:physics:phyx-mini-0779:phyxminiproblems-problemphyxmini0779-rotationsense, def:physics:phyx-mini-0779:phyxminiproblems-problemphyxmini0779-suppliedhoopfigure}
  Independent physical quantities and endpoint observables. In particular, finalCenterSpeed is not defined from an answer choice or from a solved formula; it is constrained only by the governing laws below
\end{definition}
\begin{proof}
  This is the declared model definition of Unwinding Hoop Setup.
\end{proof}
\begin{definition}[Matches Problem Scenario]
  \label{def:physics:phyx-mini-0779:phyxminiproblems-problemphyxmini0779-matchesproblemscenario}
  \lean{PhyXMiniProblems.ProblemPhyXMini0779.MatchesProblemScenario}
  \uses{def:physics:phyx-mini-0779:phyxminiproblems-problemphyxmini0779-unwindinghoopsetup}
  The qualitative physical model stated or conventionally idealized by the exercise. "Several times" is recorded conservatively as at least two full wraps; the exact count does not enter the endpoint energy calculation
\end{definition}
\begin{proof}
  This is the declared model definition of Matches Problem Scenario.
\end{proof}
\begin{definition}[Matches Problem Readouts]
  \label{def:physics:phyx-mini-0779:phyxminiproblems-problemphyxmini0779-matchesproblemreadouts}
  \lean{PhyXMiniProblems.ProblemPhyXMini0779.MatchesProblemReadouts}
  \uses{def:physics:phyx-mini-0779:phyxminiproblems-problemphyxmini0779-speedreadout, def:physics:phyx-mini-0779:phyxminiproblems-problemphyxmini0779-angularspeedreadout, def:physics:phyx-mini-0779:phyxminiproblems-problemphyxmini0779-massinkilograms, def:physics:phyx-mini-0779:phyxminiproblems-problemphyxmini0779-lengthincentimeters, def:physics:phyx-mini-0779:phyxminiproblems-problemphyxmini0779-unwindinghoopsetup}
  Numerical data from the prose and the release-from-rest initial condition. No final speed or answer choice occurs in this structure
\end{definition}
\begin{proof}
  This is the declared model definition of Matches Problem Readouts.
\end{proof}
\begin{definition}[Matches Primary Figure]
  \label{def:physics:phyx-mini-0779:phyxminiproblems-problemphyxmini0779-matchesprimaryfigure}
  \lean{PhyXMiniProblems.ProblemPhyXMini0779.MatchesPrimaryFigure}
  \uses{def:physics:phyx-mini-0779:phyxminiproblems-problemphyxmini0779-lengthinmeters, def:physics:phyx-mini-0779:phyxminiproblems-problemphyxmini0779-figureobject, def:physics:phyx-mini-0779:phyxminiproblems-problemphyxmini0779-unwindinghoopsetup}
  Literal primary-image evidence: a hand holds a vertical string tangent to the hoop, the hoop has a clockwise curved arrow, and the radial indicator is labeled 0.0800 m. It contains no final-speed readout
\end{definition}
\begin{proof}
  This is the declared model definition of Matches Primary Figure.
\end{proof}
\begin{definition}[Uses Standard Earth Gravity]
  \label{def:physics:phyx-mini-0779:phyxminiproblems-problemphyxmini0779-usesstandardearthgravity}
  \lean{PhyXMiniProblems.ProblemPhyXMini0779.UsesStandardEarthGravity}
  \uses{def:physics:phyx-mini-0779:phyxminiproblems-problemphyxmini0779-accelerationinmeterspersecondsquared, def:physics:phyx-mini-0779:phyxminiproblems-problemphyxmini0779-unwindinghoopsetup}
  Standard near-Earth gravitational acceleration used by the numerical multiple-choice calculation
\end{definition}
\begin{proof}
  This is the declared model definition of Uses Standard Earth Gravity.
\end{proof}
\begin{definition}[Has Physical Parameters]
  \label{def:physics:phyx-mini-0779:phyxminiproblems-problemphyxmini0779-hasphysicalparameters}
  \lean{PhyXMiniProblems.ProblemPhyXMini0779.HasPhysicalParameters}
  \uses{def:physics:phyx-mini-0779:phyxminiproblems-problemphyxmini0779-massreadout, def:physics:phyx-mini-0779:phyxminiproblems-problemphyxmini0779-lengthreadout, def:physics:phyx-mini-0779:phyxminiproblems-problemphyxmini0779-accelerationreadout, def:physics:phyx-mini-0779:phyxminiproblems-problemphyxmini0779-momentofinertiareadout, def:physics:phyx-mini-0779:phyxminiproblems-problemphyxmini0779-unwindinghoopsetup}
  Positivity and nondegeneracy of the physical parameters
\end{definition}
\begin{proof}
  This is the declared model definition of Has Physical Parameters.
\end{proof}
\begin{definition}[Satisfies Thin Hoop Inertia Law]
  \label{def:physics:phyx-mini-0779:phyxminiproblems-problemphyxmini0779-satisfiesthinhoopinertialaw}
  \lean{PhyXMiniProblems.ProblemPhyXMini0779.SatisfiesThinHoopInertiaLaw}
  \uses{def:physics:phyx-mini-0779:phyxminiproblems-problemphyxmini0779-massreadout, def:physics:phyx-mini-0779:phyxminiproblems-problemphyxmini0779-lengthreadout, def:physics:phyx-mini-0779:phyxminiproblems-problemphyxmini0779-momentofinertiareadout, def:physics:phyx-mini-0779:phyxminiproblems-problemphyxmini0779-unwindinghoopsetup}
  Axial moment of inertia of a uniform thin hoop, I = m R². This is a general constitutive law for the selected body model, not the requested speed
\end{definition}
\begin{proof}
  This is the declared model definition of Satisfies Thin Hoop Inertia Law.
\end{proof}
\begin{definition}[Satisfies Fixed String No Slip Kinematics]
  \label{def:physics:phyx-mini-0779:phyxminiproblems-problemphyxmini0779-satisfiesfixedstringnoslipkinematics}
  \lean{PhyXMiniProblems.ProblemPhyXMini0779.SatisfiesFixedStringNoSlipKinematics}
  \uses{def:physics:phyx-mini-0779:phyxminiproblems-problemphyxmini0779-lengthreadout, def:physics:phyx-mini-0779:phyxminiproblems-problemphyxmini0779-speedreadout, def:physics:phyx-mini-0779:phyxminiproblems-problemphyxmini0779-angularspeedreadout, def:physics:phyx-mini-0779:phyxminiproblems-problemphyxmini0779-unwindinghoopsetup}
  For a fixed string that unwinds without sliding, the center speed equals the rim radius times the angular-speed magnitude, v = R ω, at each endpoint
\end{definition}
\begin{proof}
  This is the declared model definition of Satisfies Fixed String No Slip Kinematics.
\end{proof}
\begin{definition}[Conserves Mechanical Energy During Descent]
  \label{def:physics:phyx-mini-0779:phyxminiproblems-problemphyxmini0779-conservesmechanicalenergyduringdescent}
  \lean{PhyXMiniProblems.ProblemPhyXMini0779.ConservesMechanicalEnergyDuringDescent}
  \uses{def:physics:phyx-mini-0779:phyxminiproblems-problemphyxmini0779-massreadout, def:physics:phyx-mini-0779:phyxminiproblems-problemphyxmini0779-lengthreadout, def:physics:phyx-mini-0779:phyxminiproblems-problemphyxmini0779-speedreadout, def:physics:phyx-mini-0779:phyxminiproblems-problemphyxmini0779-accelerationreadout, def:physics:phyx-mini-0779:phyxminiproblems-problemphyxmini0779-angularspeedreadout, def:physics:phyx-mini-0779:phyxminiproblems-problemphyxmini0779-momentofinertiareadout, def:physics:phyx-mini-0779:phyxminiproblems-problemphyxmini0779-unwindinghoopsetup}
  Mechanical-energy balance between release and the endpoint after descent: m g h + 1/2 m v₀² + 1/2 I ω₀² = 1/2 m v² + 1/2 I ω². The ideal string and static hand contact do no net work on the combined hoop motion in this model. The law is stated in every coherent system of selected mechanical base units and contains no solved final-speed value
\end{definition}
\begin{proof}
  This is the declared model definition of Conserves Mechanical Energy During Descent.
\end{proof}
\begin{lemma}[final Center Speed eq sqrt gravity mul descent]
  \label{lem:physics:phyx-mini-0779:phyxminiproblems-problemphyxmini0779-finalcenterspeed-eq-sqrt-gravity-mul-descent}
  \lean{PhyXMiniProblems.ProblemPhyXMini0779.finalCenterSpeed_eq_sqrt_gravity_mul_descent}
  \uses{def:physics:phyx-mini-0779:phyxminiproblems-problemphyxmini0779-lengthinmeters, def:physics:phyx-mini-0779:phyxminiproblems-problemphyxmini0779-speedinmeterspersecond, def:physics:phyx-mini-0779:phyxminiproblems-problemphyxmini0779-accelerationinmeterspersecondsquared, def:physics:phyx-mini-0779:phyxminiproblems-problemphyxmini0779-unwindinghoopsetup, def:physics:phyx-mini-0779:phyxminiproblems-problemphyxmini0779-matchesproblemreadouts, def:physics:phyx-mini-0779:phyxminiproblems-problemphyxmini0779-hasphysicalparameters, def:physics:phyx-mini-0779:phyxminiproblems-problemphyxmini0779-satisfiesthinhoopinertialaw, def:physics:phyx-mini-0779:phyxminiproblems-problemphyxmini0779-satisfiesfixedstringnoslipkinematics, def:physics:phyx-mini-0779:phyxminiproblems-problemphyxmini0779-conservesmechanicalenergyduringdescent}
  Combining release from rest, energy conservation, I = m R², and v = R ω makes the equal translational and rotational kinetic-energy terms sum to m v². Hence the nonnegative center speed is sqrt (g h)
\end{lemma}
\begin{proof}
  The conclusion follows from the governing relations and figure readouts named in the statement of final Center Speed eq sqrt gravity mul descent.
\end{proof}
\begin{definition}[Answer Choice]
  \label{def:physics:phyx-mini-0779:phyxminiproblems-problemphyxmini0779-answerchoice}
  \lean{PhyXMiniProblems.ProblemPhyXMini0779.AnswerChoice}
  Labels of the four speed choices printed with the problem
\end{definition}
\begin{proof}
  This is the declared model definition of Answer Choice.
\end{proof}
\begin{definition}[displayed Speed In Meters Per Second]
  \label{def:physics:phyx-mini-0779:phyxminiproblems-problemphyxmini0779-displayedspeedinmeterspersecond}
  \lean{PhyXMiniProblems.ProblemPhyXMini0779.displayedSpeedInMetersPerSecond}
  \uses{def:physics:phyx-mini-0779:phyxminiproblems-problemphyxmini0779-answerchoice}
  Printed speed beside each answer choice, in meters per second
\end{definition}
\begin{proof}
  This is the declared model definition of displayed Speed In Meters Per Second.
\end{proof}
\begin{definition}[recorded Dataset Answer]
  \label{def:physics:phyx-mini-0779:phyxminiproblems-problemphyxmini0779-recordeddatasetanswer}
  \lean{PhyXMiniProblems.ProblemPhyXMini0779.recordedDatasetAnswer}
  \uses{def:physics:phyx-mini-0779:phyxminiproblems-problemphyxmini0779-answerchoice}
  Dataset metadata; this declaration is not used as a theorem premise
\end{definition}
\begin{proof}
  This is the declared model definition of recorded Dataset Answer.
\end{proof}
\begin{definition}[Rounds To Displayed Hundredth]
  \label{def:physics:phyx-mini-0779:phyxminiproblems-problemphyxmini0779-roundstodisplayedhundredth}
  \lean{PhyXMiniProblems.ProblemPhyXMini0779.RoundsToDisplayedHundredth}
  \uses{def:physics:phyx-mini-0779:phyxminiproblems-problemphyxmini0779-speedinmeterspersecond, def:physics:phyx-mini-0779:phyxminiproblems-problemphyxmini0779-unwindinghoopsetup, def:physics:phyx-mini-0779:phyxminiproblems-problemphyxmini0779-answerchoice, def:physics:phyx-mini-0779:phyxminiproblems-problemphyxmini0779-displayedspeedinmeterspersecond}
  The exact modeled speed rounds to a displayed hundredth-meter-per-second value
\end{definition}
\begin{proof}
  This is the declared model definition of Rounds To Displayed Hundredth.
\end{proof}
\begin{definition}[Is Unique Closest Displayed Speed]
  \label{def:physics:phyx-mini-0779:phyxminiproblems-problemphyxmini0779-isuniqueclosestdisplayedspeed}
  \lean{PhyXMiniProblems.ProblemPhyXMini0779.IsUniqueClosestDisplayedSpeed}
  \uses{def:physics:phyx-mini-0779:phyxminiproblems-problemphyxmini0779-speedinmeterspersecond, def:physics:phyx-mini-0779:phyxminiproblems-problemphyxmini0779-unwindinghoopsetup, def:physics:phyx-mini-0779:phyxminiproblems-problemphyxmini0779-answerchoice, def:physics:phyx-mini-0779:phyxminiproblems-problemphyxmini0779-displayedspeedinmeterspersecond}
  The physical speed is strictly closer to one displayed value than to every other displayed value
\end{definition}
\begin{proof}
  This is the declared model definition of Is Unique Closest Displayed Speed.
\end{proof}
% --- Archon declaration topology end ---
% --- Archon physics formalization source end ---
